\chapter{Future Work}\label{ch:futurework}

\textcolor{red}{AntCalc}, as described in Chapter~\ref{ch:packageframework}
and validated in Chapter~\ref{ch:validation}, for both its the building and integration 
engines, currently encompasses all massless antennae in SMQCD up to NNLO.

Modern fixed-order QCD calculations, however, increasingly require antenna functions beyond this
scope. Massive quark effects enter heavy-flavour production at hadron colliders, supersymmetric
models introduce new parton species, and Higgs-plus-jet processes at NNLO require antennae
derived from an effective $H\to gg$ vertex. Extension to N$^3$LO is also an active frontier,
as the precision demands of the LHC physics programme continue to tighten. The natural 
directions for extending \textcolor{red}{AntCalc} are therefore the massive SMQCD sector,
the massless SUSY and HiggsEFT models, and eventually the full N$^3$LO antenna set.

The physical, mathematical and computation challenges introduced by these extensions are discussed
to a higher level of detail in Appendix~\ref{ap:futureWork}.

\section{Massive Antennae}

As discussed in Chapter~\ref{ch:workedexample}, the $A_3^0$ antenna is the simplest non-trivial
SMQCD antenna in the model, and as such, is also the best topology to use to extend 
\textcolor{red}{AntCalc} to the massive framework. This extension has already been 
implemented in the current version. To access the results, we can run the usual functions, 
using the \texttt{quarkMass} option,
\begin{mdframed}[backgroundcolor=gray!10, linewidth=0pt]
    \texttt{In[2]:= BuildAntenna[A,3,0,quarkMass->mQ]}\\
    \hfill
    \texttt{Out[2]:=}
        \begin{equation*}
            \begin{split}
                \frac{1}{4s_{13}^2s_{23}^2\left(4m_q^2+s_{12}+s_{13}+s_{23}\right)} \Bigg[&-8m_q^2\left(s_{13}^2+s_{23}^2\right)+s_{13}s_{23}\left(2s_{12}^2+s_{13}^2+s_{23}^2+2s_{12}(s_{13}+s_{23})\right)\\
                                                                                          &-2m_q^2\Big(s_{13}^2+s_{13}^2s_{23}+s_{13}s_{23}^2+s_{23}^2+s_{12}\left(s_{13}^2-4s_{13}s_{23}+s_{23}^2\right)\Big)\Bigg]
            \end{split}
        \end{equation*}
    \texttt{In[3]:= \\
            BuildAndIntegrateAntenna[A,3,0,quarkMass->mQ]\\
            /.\{j[MX30Basis123, 1, 1, 1, 0, 0]->RM1, \\
            j[MX30Basis123, 2, 1, 1, 0, 0]->RM2,d->4-2Epsilon, eps->Epsilon\}\\
            //FullSimplify//Collect[\#,\{RM1,RM2\}]\&}\\
    \hfill
    \texttt{Out[3]:=}
        \begin{equation*}
            \begin{split}
                &\frac{8\big(2-\epsilon(7-6\epsilon)\big)m_q^4-4\Big(4+\epsilon\big(-11+4\epsilon(1+\epsilon)\big)\Big)m_q^2q^2-(1-2\epsilon)\big(4-\epsilon(3+2\epsilon)\big)q^4}{\epsilon(1+2\epsilon)\left(4m_q^2-q^2\right)q^4}R_1^M\\
                &+\frac{16(1-2\epsilon)m_q^6+8\big(-1+4\epsilon(1+\epsilon)\big)m_Q^4q^2-4\left(1+2\epsilon+4\epsilon^2\right)m_q^2q^4+2q^6+\epsilon(-1+2\epsilon)q^6}{\epsilon(1-2\epsilon)\left(4m_q^2-q^2\right)q^4}R_2^M
            \end{split}
        \end{equation*}
\end{mdframed}

To get to these results, the same method from Chapter~\ref{ch:workedexample} was used, but
for the massive case, in which $k^2\neq0\neq m_q$. The results still are not complete, as 
the massive master integrals, $R_1^M$ and $R_2^M$, have yet to be computed to enable us to 
expand into the usual $\epsilon$ Laurent series.
This example does show, however, that the already present framework is able to handle the 
massive expansion, at least inside SMQCD, seamlessly, opening it up to the possibility of 
further expansion to the massive SMQCD (and beyond) sector.

\section{SUSY model}

\subsection{Short Introduction To SUSY}

\subsection{Example Implementation}

\section{HiggsEFT model}

\subsection{Short Introduction To HiggsEFT}

\subsection{Example Implementation}

\section{N$^3$LO}